\chapter*{Геймификаци, тоглоомд суурилсан суралцахуй ба оролцооны онолын иж бүрэн шинжилгээ}
\begin{sloppy}

\section*{Гүйцэтгэх хураангуй}
Геймификаци нь тоглоомын дизайны элементүүдийг тоглоом биш орчинд ашиглахыг хэлдэг бол тоглоомд суурилсан суралцахуй нь өөрөө тоглоомын бүтэц, дүрэм, хийсвэр зөрчил, сорилт, хариу үйлдэл, буцах холбоотой орчныг суралцах зорилгод зориулан зохион бүтээдэг.

Эмпирик судалгааны том зураг нь ``геймификаци үр дүнтэй байж чадна, гэхдээ автоматаар биш'' гэсэн дүгнэлт рүү хүрдэг. Тэр дундаа контекст, хэрэглэгчийн төрөл, механикийн хослол, нийгмийн горим зэргээс үр нөлөө нь ихээхэн хамаардаг. Чадвар–сорилтын тэнцвэрийг хангах үед үр нөлөө өсдөг боловч ``оноо, тэмдэг, зэрэглэл = мотиваци'' гэсэн шулуун шугаман ойлголт практик дээр учир дутагдалтай болохыг судалгаанууд харуулсаар байна.

\section*{Үндсэн ойлголт ба түүхэн хөгжил}
Геймификацийн орчин үеийн академик тодорхойлолтыг Sebastian Deterding 2011 онд ``game design elements''-ийг ``non-game contexts''-д хэрэглэх гэж томьёолсон. Тоглоомд суурилсан суралцахуйг (Game-Based Learning) Plass, Homer нар геймификациас ялгаж, challenge–response–feedback гэсэн циклийг сургалтын үндсэн нэгж гэж үзсэн.

Түүхэн хувьд энэ салбар нь дараах үе шатуудыг дамжин хөгжжээ:
\begin{itemize}
    \item 1957: Reinforcement schedules (B.F. Skinner)
    \item 1977: Self-efficacy (Bandura)
    \item 1990: Flow (Mihaly Csikszentmihalyi)
    \item 2004: MDA framework, Immersion, School engagement
    \item 2011: Gamification definition (Deterding et al.)
    \item 2020: Meta-analyses of gamified learning
\end{itemize}

\section*{Гол онолууд ба загварууд}
Сайн дизайн нь зөвхөн нэг онол бус, дараах концепциудыг хослуулан уншиж байж оновчтой болдог:
\begin{itemize}
    \item \textbf{Өөрөө тодорхойлох онол (Self-Determination Theory):} Autonomy, Competence, Relatedness гэсэн гурван суурь хэрэгцээг хангах.
    \item \textbf{Зорилго тавих онол (Goal-Setting Theory):} Тодорхой, сорилттой зорилго нь гүйцэтгэлийг чиглүүлж, хүчин чармайлтыг өсгөнө.
    \item \textbf{Урсгалын онол (Flow Theory):} Сорилт ба чадварын тэнцвэр, тодорхой зорилго, шууд хариу үйлдэл (feedback).
    \item \textbf{MDA Framework:} Mechanics (Механик), Dynamics (Динамик), Aesthetics (Эстетик) гэсэн гурван түвшний уялдаа.
\end{itemize}

\section*{Механизм ба Дизайн паттерн}
Геймификацийн механик нь динамикаар дамжин сэтгэлзүйн хэрэгцээг хангаж, улмаар зан үйлийн үр дүнд хүргэдэг.
\begin{itemize}
    \item \textbf{Оноо (Points):} Үйлдлийг тоон хэлбэрт оруулж, ахиц дэвшлийг харагдуулах (Micro-reward).
    \item \textbf{Түвшин (Levels):} Урт хугацааны progression ба статусыг тодорхойлох.
    \item \textbf{Тэмдэг (Achievements/Badges):} Амжилтыг тэмдэглэх, ур чадварыг баталгаажуулах.
    \item \textbf{Сорилт (Challenge/Quest):} Хэрэглэгчийн үйлдлийг сонирхолтой болгох, Flow төлөвт оруулах.
    \item \textbf{Зэрэглэл (Leaderboard):} Нийгмийн харьцуулалт ба өрсөлдөөнийг өдөөх.
    \item \textbf{Зуршил (Streak):} Давтамжтай зуршил үүсгэх, Loss aversion механизмыг ашиглах.
\end{itemize}

\section*{Оролцоо ба Immersion-ийг хэмжих үзүүлэлтүүд}
Геймификацийн хэмжилтэд зөвхөн ганц KPI (жишээ нь, системд зарцуулсан хугацаа) хангалтгүй. Дараах хэмжигдэхүүнүүдийг ашиглахыг зөвлөдөг:
\begin{itemize}
    \item \textbf{User Engagement Scale (UES):} Гоо зүйн таашаал, төвлөрөл, шинэлэг байдал, хэрэглэгдэхүүн байдал.
    \item \textbf{Flow State Scale:} Сэтгэл ханамж, цаг хугацааны мэдрэмж, хяналт.
    \item \textbf{Behavioral Metrics:} Retention, dwell time, task success rate, repeat behavior.
\end{itemize}

Дүгнэж хэлэхэд, геймификаци нь зөвхөн ``тоглоом болгох'' бус, харин хэрэглэгчийн дотоод сэдэлжүүлэлтийг (intrinsic motivation) дэмжсэн утга учиртай орчныг бүрдүүлэх тухай асуудал юм.

\end{sloppy}
